\section{统计实操补充：如何设定 ``可发布'' 的显著性门槛}

\subsection{从 effect size 到 MDE（Minimum Detectable Effect）}
课程反复强调 ``不要把微小分差直接解读成能力提升''。在工程实践中，可把这个原则落实为 MDE 机制：在发布前先声明 ``至少提升多少才算有效进步''，低于该阈值则进入观察区而非发布区。

\begin{knowledgebox}{MDE 的直观解释}
MDE 不是统计上的真理，而是团队关于 ``业务上有意义提升'' 的共同约定。它把统计显著性和产品价值连接起来，避免为了追逐小数点后两位而过度优化 benchmark。
\end{knowledgebox}

若把每次运行结果近似视为独立样本，均值比较的样本量可用如下近似式估计：
\[
n \approx 2 \times \left(\frac{z_{1-\alpha/2}+z_{1-\beta}}{\Delta/\sigma}\right)^2
\]
其中 $\Delta$ 为目标检测差异（MDE），$\sigma$ 为历史波动标准差，$\alpha$ 为显著性水平，$\beta$ 为二类错误率。该式给出的不是绝对答案，但可快速评估 ``当前 run 数是否明显不足''。

\begin{warningbox}{常见误区：先跑完再找显著性}
很多团队先跑几次，看到结果 ``似乎更好'' 才临时做显著性检验。这样会引入选择偏差。更稳健做法是提前写下评测计划：run 次数、停止条件、主指标与次指标、判定阈值。
\end{warningbox}

\subsection{重复运行预算：什么时候 5 次够，什么时候远远不够}
对于低噪声分类任务，5 次重复可能已经足够稳定；但对 agent benchmark，轨迹随机性和环境扰动会使方差显著上升，10--20 次运行才可能得到可解释区间。课程中的建议是基于历史方差做分层预算，而不是给所有任务统一 run 数。

\begin{table}[H]
\centering
\begin{tabular}{p{3.0cm}p{2.8cm}p{3.4cm}p{2.9cm}}
\toprule
\textbf{任务类型} & \textbf{历史波动（示例）} & \textbf{建议 run 数} & \textbf{发布建议} \\
\midrule
Closed QA / MCQ & $\sigma \le 0.3$ & 5--8 & 报告均值+CI，通常可快速发布 \\
Open-ended generation & $0.3 < \sigma \le 1.0$ & 8--12 & 强制双 judge 抽检 \\
Agent workflow & $\sigma > 1.0$ & 12--20 & 必报轨迹级错误与环境版本 \\
\bottomrule
\end{tabular}
\caption{按噪声强度配置运行预算，比固定 run policy 更稳健}
\end{table}

\begin{importantbox}{可执行规则}
当观测提升 $< \mathrm{MDE}$ 或 95\%CI 与基线显著重叠时，默认结论应为 ``暂无证据表明显著提升''，而非 ``模型退化/提升'' 的二元叙事。
\end{importantbox}

\subsection{多指标冲突：如何从 ``单分数'' 过渡到 ``发布门禁''}
真实系统几乎总会出现指标冲突，例如 success rate 上升但 timeout 上升，或总体分数上升但 hard split 下降。课程建议把发布判定写成门禁规则（gating policy），例如：
\begin{itemize}
    \item 主指标达到 MDE 且通过显著性门槛；
    \item 可靠性指标（timeout、invalid action）不得劣化超过阈值；
    \item hard split 不得出现统计显著退化；
    \item 若使用 LLM-as-judge，必须通过抽样人工复核。
\end{itemize}

\begin{knowledgebox}{为什么要 ``门禁化''}
门禁化的核心价值是降低解释自由度。发布前就定义好规则，可以防止事后挑选有利指标，也让跨版本比较具有一致标准。
\end{knowledgebox}

\subsection{本章小结}
统计实操的关键不是追求复杂检验，而是把 ``阈值、预算、门禁'' 前置定义。这样 benchmark 才能从研究展示工具升级为工程决策基础。

